\section{Results}
\label{sec:results}
\textcolor{red}{\textbf{Evidence in this version: the synthetic check of the statistic (Section~\ref{sec:toy}) and a supervised first-stage run to step 2{,}000 (Section~\ref{sec:stage1}, three held-out views of one scene only). No guided sampling, semi-self-supervised training, calibration or fault-injection experiment has been run; any other table or figure is a placeholder.}}

\subsection{Synthetic check of the test statistic (toy, not real data)}
\label{sec:toy}
The statistic of Section~\ref{sec:statistic} was exercised on a synthetic Gaussian world (a $16\times16$ grid, $M=10$ samples per view, one seed,
$200$ clean trials for the null and $80$ trials per fault cell, power standard error about $0.05$); the render, the unguided samples and the truth are
drawn from known Gaussian processes, the truth being one more draw of the posterior. No image and no learned model is involved; the purpose is to check the
statistic and its calibration, not to estimate the performance of the method. Five variants are compared: the true scale and shift with the true covariance
(unattainable), the fitted scale and shift with the true covariance (oracle), and the fitted scale and shift with a stationary-kernel covariance estimated from the
samples (robust and plain fit), and with a shrunk sample covariance. The threshold is the empirical 95\% quantile of the clean-view null of the identical pipeline.

\begin{figure}[!t]
  \centering
  \setlength{\figw}{\linewidth}\adjustbox{max width=\linewidth}{% Toy statistic: power versus fault magnitude. Data: data/fig3_{square,warp,edge}.csv (synthetic Gaussian toy, 1 seed, 80 trials/cell)
\newcommand{\figthreeseries}[1]{%
  \addplot[dotted, black, line width=0.5pt, forget plot] coordinates {(0.3,0.8) (4.2,0.8)};
  \addplot[dotted, black, line width=0.5pt, forget plot] coordinates {(0.3,0.05) (4.2,0.05)};
  \figthreeone{known}{known}{#1}
  \figthreeone{oracle}{oracle}{#1}
  \figthreeone{kernel}{kernel}{#1}
  \figthreeone{kernel_nr}{kernelnr}{#1}
  \figthreeone{shrunk}{shrunk}{#1}}
\newcommand{\figthreeone}[3]{\addplot[#2, error bars/.cd, y dir=both, y explicit, error bar style={line width=0.4pt},
  error mark options={rotate=90, mark size=1.0pt, line width=0.4pt}]
  table[x=kappa, y=#1, y error=#1_se, col sep=comma]{\tikzdir/data/fig3_#3.csv};}
\begin{tikzpicture}
\begin{groupplot}[amb, group style={group size=3 by 1, horizontal sep=0.075\figw, y descriptions at=edge left},
  width=0.235\figw, height=0.2\figw, xmin=0.3, xmax=4.2, ymin=0, ymax=1.08, ytick={0,0.2,0.4,0.6,0.8,1.0}, xtick={1,2,3,4},
  xlabel={Fault magnitude [sd]}]
\nextgroupplot[ylabel={Power [-]}, legend to name=figthreeleg, legend columns=2,
  legend style={/tikz/every even column/.append style={column sep=6pt}}]
\figthreeseries{square}
\addlegendentry{known: true $(a,b)$ and covariance (unattainable)}
\addlegendentry{fitted $(s,t)$, true covariance}
\addlegendentry{fitted $(s,t)$, estimated kernel (robust)}
\addlegendentry{fitted $(s,t)$, estimated kernel (plain)}
\addlegendentry{fitted $(s,t)$, shrunk sample covariance}
\nextgroupplot
\figthreeseries{warp}
\nextgroupplot
\figthreeseries{edge}
\end{groupplot}
\node[anchor=south] at ($(group c2r1.north)+(0,0.25cm)$) {\ref{figthreeleg}};
\subcap{$(group c1r1.south)+(0,-0.95cm)$}{(a) Square region}
\subcap{$(group c2r1.south)+(0,-0.95cm)$}{(b) Smooth warp}
\subcap{$(group c3r1.south)+(0,-0.95cm)$}{(c) Edge-aligned region}
\end{tikzpicture}
}
  \caption{Power against fault magnitude (units of the per-pixel total residual standard deviation) for three fault families. \emph{Synthetic Gaussian toy, one seed,
  80 trials per cell; error bars are binomial standard errors.} Dotted lines mark 80\% power and the 5\% level.}
  \label{fig:toy-power}
\end{figure}

With the true covariance, fitting the scale and shift costs almost nothing for local faults: the minimum detectable fault at 80\% power is $0.88$ standard
deviations for a square region with the fitted scale and shift against $0.89$ with the true ones, and $0.84$ against $0.89$ for an edge-aligned region; a smooth
warp needs about $2.1$ (Figure~\ref{fig:toy-power}, \ref{fig:toy-mdf}). The covariance estimate dominates the loss: with the estimated kernel the same faults need
$2.2$--$2.4$ and the warp is not detected up to $4$ standard deviations. A constant offset is invisible by construction (power $0.05$), whereas a scale error is
flagged (power $1.0$ with the true covariance), because the fitted scale also rescales the sample covariance. With the $\chi^2_{n-2}$ reference the clean-view
false-alarm rate is $0.06$ with the true covariance but $0.27$ with the estimated kernel, which is why the threshold is taken from the empirical null.

\begin{figure}[!t]
  \centering
  \setlength{\figw}{\linewidth}\adjustbox{max width=\linewidth}{% Toy statistic: minimum detectable fault (80 % power) per covariance variant, and false alarms at the chi2 threshold.
% Data: data/fig4_mdf_wide.csv, data/fig4_null.csv (synthetic Gaussian toy, 1 seed)
\newcommand{\figfourbar}[4]{% #1 column, #2 colour/style, #3 shift, #4 legend text
  \addplot[ybar, #2, draw=black, line width=0.3pt, bar shift=#3, bar width=7pt, nodes near coords, point meta=explicit symbolic,
    every node near coord/.append style={font=\figfont\tiny, rotate=90, anchor=west, /pgf/number format/.cd}]
    table[x expr=\coordindex, y=#1, meta=#1_lab, col sep=comma]{\tikzdir/data/fig4_mdf_wide.csv};
  \addplot[ybar, forget plot, fill=none, draw=none, bar shift=#3, bar width=7pt, postaction={pattern=north east lines, pattern color=white}]
    table[x expr=\coordindex, y expr=\thisrow{#1_c}*\thisrow{#1}, col sep=comma]{\tikzdir/data/fig4_mdf_wide.csv};
  \addlegendentry{#4}}
\pgfplotsset{figfour/.style={ybar, barleg}}
\newcommand{\figfournull}[2]{\addplot[ybar, fill=#2, draw=black, line width=0.3pt, bar shift=0pt, bar width=9pt, nodes near coords,
  restrict expr to domain={\thisrow{vi}}{#1:#1}, every node near coord/.append style={font=\figfont\tiny, /pgf/number format/.cd, fixed, precision=3}]
  table[x=vi, y=far, col sep=comma]{\tikzdir/data/fig4_null.csv};}
\begin{tikzpicture}
\begin{groupplot}[amb, group style={group size=2 by 1, horizontal sep=0.1\figw}, height=0.22\figw]
\nextgroupplot[figfour, width=0.43\figw, ymin=0, ymax=5.2, ytick={0,1,2,3,4}, xtick={0,1,2}, xticklabels={Square, Smooth warp, Edge-aligned}, xmin=-0.55, xmax=2.55,
  ylabel={Minimum detectable fault [sd]}, legend to name=figfourleg, legend columns=2, legend style={/tikz/every even column/.append style={column sep=6pt}}]
\figfourbar{known}{fill=ambBlack}{-14pt}{known: true $(a,b)$ and covariance}
\figfourbar{oracle}{fill=ambRed}{-7pt}{fitted $(s,t)$, true covariance}
\figfourbar{kernel}{fill=ambBlue}{0pt}{estimated kernel (robust)}
\figfourbar{kernel_nr}{fill=ambBrown}{7pt}{estimated kernel (plain)}
\figfourbar{shrunk}{fill=ambGreen}{14pt}{shrunk sample covariance}
\nextgroupplot[figfour, width=0.25\figw, ymin=0, ymax=0.34, xtick={0,1,2,3,4}, xticklabels={known, oracle, kernel, kernel (plain), shrunk}, x tick label style={rotate=30, anchor=east},
  xmin=-0.6, xmax=4.6, ylabel={False-alarm rate [-]}]
\draw[dotted, black, line width=0.5pt] (axis cs:-0.6,0.05) -- (axis cs:4.6,0.05);
\figfournull{0}{ambBlack}\figfournull{1}{ambRed}\figfournull{2}{ambBlue}\figfournull{3}{ambBrown}\figfournull{4}{ambGreen}
\end{groupplot}
\node[anchor=south] at ($(group c1r1.north east)+(0.8cm,0.3cm)$) {\ref{figfourleg}};
\subcap{$(group c1r1.south)+(0,-1.0cm)$}{(a) Detectability (hatched: not reached)}
\subcap{$(group c2r1.south)+(0,-1.35cm)$}{(b) False alarms at $\chi^2$ threshold}
\end{tikzpicture}
}
  \caption{Minimum detectable fault at 80\% power per covariance variant (hatched: not reached by $4$) and clean-view false-alarm rate at the
  $\chi^2_{n-2}$ 95\% threshold. \emph{Synthetic Gaussian toy, one seed, 200 clean trials and 80 trials per cell.}}
  \label{fig:toy-mdf}
\end{figure}

\begin{figure}[!t]
  \centering
  \setlength{\figw}{\linewidth}\adjustbox{max width=\linewidth}{% Toy statistic under a miscalibrated sample spread. Data: data/fig6_{oracle,kernel,kernel_nr}.csv (synthetic Gaussian toy, 1 seed)
\newcommand{\figsixline}[3]{% #1 file, #2 style, #3 y column
  \addplot[#2] table[x=miscal, y=#3, col sep=comma]{\tikzdir/data/fig6_#1.csv};}
\newcommand{\figsixcens}[3]{% censored points: hollow marker and ">4"
  \addplot[only marks, mark=o, mark size=3pt, #2, nodes near coords={$>4$}, every node near coord/.append style={font=\figfont\tiny, anchor=south},
    restrict expr to domain={\thisrow{#3}}{1:1}] table[x=miscal, y=mdf_warp, col sep=comma]{\tikzdir/data/fig6_#1.csv};}
\begin{tikzpicture}
\begin{groupplot}[amb, group style={group size=3 by 1, horizontal sep=0.115\figw}, width=0.235\figw, height=0.2\figw,
  xmode=log, xmin=0.5, xmax=1.9, xtick={0.6,1,1.6}, xticklabels={0.6,1.0,1.6}, minor tick num=0, xlabel={Spread / true spread [-]}]
\nextgroupplot[ylabel={Mean $T/(n-2)$ [-]}, ymin=0.2, ymax=2.9, legend to name=figsixleg, legend columns=3,
  legend style={/tikz/every even column/.append style={column sep=6pt}}]
\addplot[dotted, black, line width=0.5pt, forget plot] coordinates {(0.5,1) (1.9,1)};
\figsixline{oracle}{oracle}{t_over_dof}\figsixline{kernel}{kernel}{t_over_dof}\figsixline{kernel_nr}{kernelnr}{t_over_dof}
\addlegendentry{true covariance}\addlegendentry{est. kernel (robust)}\addlegendentry{est. kernel (plain)}
\nextgroupplot[ylabel={MDF [sd]}, ymin=0.6, ymax=2.6]
\figsixline{oracle}{oracle}{mdf_sq}\figsixline{kernel}{kernel}{mdf_sq}\figsixline{kernel_nr}{kernelnr}{mdf_sq}
\nextgroupplot[ylabel={MDF [sd]}, ymin=1.8, ymax=4.5]
\figsixline{oracle}{oracle}{mdf_warp}\figsixline{kernel}{kernel}{mdf_warp}\figsixline{kernel_nr}{kernelnr}{mdf_warp}
\figsixcens{kernel}{ambBlue}{cens_warp}\figsixcens{kernel_nr}{ambBrown}{cens_warp}
\end{groupplot}
\node[anchor=south] at ($(group c2r1.north)+(0,0.25cm)$) {\ref{figsixleg}};
\subcap{$(group c1r1.south)+(0,-1.15cm)$}{(a) Null statistic}
\subcap{$(group c2r1.south)+(0,-1.15cm)$}{(b) Square region}
\subcap{$(group c3r1.south)+(0,-1.15cm)$}{(c) Smooth warp}
\end{tikzpicture}
}
  \caption{Effect of a miscalibrated sample spread (spread equal to $0.6$, $1.0$ and $1.6$ times the truth's posterior spread; the covariance used by the
  oracle is the one the model would report; MDF = minimum detectable fault at 80\% power, all variants with the fitted scale and shift). \emph{Synthetic Gaussian toy, one seed, 150 clean trials and 60 trials per cell.}}
  \label{fig:toy-miscal}
\end{figure}

A wrong scale of the spread is absorbed by the empirical threshold: with the oracle shape the minimum detectable fault stays at $0.87$--$0.91$ (square) and
$2.07$--$2.11$ (warp) across Figure~\ref{fig:toy-miscal}. A wrong covariance shape is what costs power. Faults covering up to $85\%$ of the view remain detectable
and collapse only at $95\%$, where the fault is nearly a constant offset; the statistic detects an inconsistency but cannot say which side is wrong.
\textcolor{red}{\textbf{Not yet tested: a calibration set that is not exchangeable with the test views, real model samples, several seeds.}}

\begin{figure}[!t]
  \centering
  \setlength{\figw}{\linewidth}\adjustbox{max width=\linewidth}{\def\toyfault{square}\input{\tikzdir/toy_render_common}
}
  \caption{One synthetic view without and with a square fault of $3$ standard deviations (the same random draw; black outline = fault support), with the true
  covariance (top) and the estimated kernel covariance (bottom). Left pair: raw residual after the fitted scale and shift. Right pair: whitened residual
  (Cholesky innovations in raster order, unit variance under the null; their streaks follow the whitening order and do not localise the fault); the test statistic
  $T$ and the empirical threshold are printed on each. \emph{Synthetic Gaussian toy, seed 7, threshold from 150 clean views.}}
  \label{fig:toy-square}
\end{figure}

\begin{figure}[!t]
  \centering
  \setlength{\figw}{\linewidth}\adjustbox{max width=\linewidth}{\def\toyfault{warp}\input{\tikzdir/toy_render_common}
}
  \caption{As Figure~\ref{fig:toy-square} for a smooth warp of $3$ standard deviations. The statistic with the true covariance flags the faulted view; with the
  estimated kernel covariance it does not ($T=412$ against a threshold of $746$), consistent with Figure~\ref{fig:toy-mdf}. \emph{Synthetic Gaussian toy, seed 7.}}
  \label{fig:toy-warp}
\end{figure}

\todofig{calibration}{Calibration of the unguided sample spread against ground truth on held-out scenes (coverage and PIT), versus the Marigold-style ensemble.}
\todofig{power}{Fault-injection power curves by fault family on real model samples, and absorption ratio of the learned prior.}

\subsection{Supervised first stage (sanity check, not a result)}
\label{sec:stage1}
\runin{Setup.} SDXL base frozen in fp16 with a fp32 VAE, LoRA rank 32 on the attention projections, a widened $8$-channel input convolution, $v$-prediction with zero-terminal-SNR, empty-prompt conditioning, $512$\,px, batch $32$.
Training data: BlendedMVS (107 training scenes, 6 held out) and a Hypersim subset; the BlendedMVS images are blended composites, not photographs.
The run was stopped at step $2{,}000$ of a planned $10{,}000$ (about $14$\,s per step on a TITAN RTX; peak $8.9$\,GB in the smoke run).
\runin{Validation.} Validation loss at steps $500/1000/1500/2000$: $0.439/0.393/0.351/0.337$; the proxy affine AbsRel (a different normalisation from the literature, a trend indicator only): $0.719/0.648/0.641/0.618$.
\runin{Held-out ETH3D \texttt{pipes}} (three views of one scene; five DDIM samples per view, scale and shift fitted to the sparse laser depth, AbsRel in depth units). Mean AbsRel over the three views: $0.102$ (step $1000$, $10$ steps), $0.105/0.107/0.110/0.111$ (step $2000$, $1/4/10/30$ steps). On view 6 the untrained base scores $0.38$.
Step $2000$ is not clearly better than step $1000$, and fewer sampler steps score slightly better; with three views of one scene these differences are within noise.
This is not a benchmark: it says only that the pipeline trains without numerical failure and moves the output toward the ground truth.
